\chapter{Introduction}\label{chapter:introduction}

This chapter sets the scene for the whole report. Background and problem come first, followed by the motivation, the objectives, a short account of the methodology, the expected outcome, and the structure of the remaining chapters.

\section{Project Overview}\label{sec:project-overview}

Large language models today have been shown to be capable of handling a wide variety of reasoning tasks such as answering scientific questions, serving as tutors or providing decision support. A large language model can still be confidently wrong. It may state an incorrect answer in fluent and assertive language that hides the error. To address this, the use of multi-agent debate has been suggested~\cite{du2024improving}. In a debate, multiple agents respond to the same question, and they review the rationales provided by other agents over several iterations. After reviewing all rationales, the agents select a single answer~\cite{du2024improving}. Chain-of-thought prompting gives each agent a way to expose its intermediate steps~\cite{wei2022chain}, which makes it easier for others to evaluate the rationale behind the agent's position. A major hurdle has been encountered at the end of this process. The major part of the problem consists of developing a procedure for combining individually generated answers produced by multiple agents. Most current systems for resolving debates using a voting mechanism~\cite{wang-etal-2024-rethinking-bounds} as the primary method for deciding which agent's response will be accepted. Each agent within such a voting system is given an equal amount of weight (i.e., each vote has an equal amount of value) and the answer that receives the most votes is chosen as the prevailing answer~\cite{choi2025debate}. Current systems determine correctness simply based upon how many times an answer appears amongst all agents. Scientific inquiry does not assess the validity of an argument by its repetition, but rather by the amount of credible evidence supporting that argument. The paper discusses a trust-calibrated multi-agent deliberation framework. This new framework will provide different levels of influence to each agent during a debate depending on how accurately they represent the information obtained via retrieving information from outside sources. Therefore, when evaluating the validity of a statement in the context of a multi-agent deliberative environment, the study suggests using a measure of evidence-grounded trust rather than a raw count of heads.

\section{Motivation}\label{sec:motivation}

Recent studies on multi-agent debate have shown that reaching consensus does not necessarily imply reaching the correct answer. Instead, debating agents may imitate one another or change their opinions under social pressure, even when one agent initially provides the correct answer. Therefore, in over 20\% of times when an agent makes a wrong decision during debate (i.e., gives a wrong answer), they ignore the correct answer made by another agent. Additionally, there is also work being done regarding modeling behavior that suggests similar tendencies. Models that are trained to be agreeable become more sycophantic and inaccurate ~\cite{ibrahim2026training}. Majority voting is particularly problematic with regard to these issues because a large number of confident voters can cause a minority who believe in the correct answer to stop believing in their answer and vote incorrectly. In addition to producing an incorrect result, majority voting creates the illusion of consensus behind which lies the actual incorrectness. These problems are most serious for scientific questions, as users may trust what they are told when they should be looking more closely at the information presented. Each existing approach addresses only a portion of the problems discussed. For example, iMAD determines if a debate is worth conducting prior to doing so ~\cite{fan2026imad}, DebUnc increases the influence of agents based upon self-reported confidence levels ~\cite{yoffe-etal-2025-debunc}, and Mixture-of-Agents combines the output from multiple agents using static aggregation ~\cite{wang2025mixture}. However, none of these approaches tie an agent's level of influence within a debate to whether that agent's claims continue to hold water when compared to external evidence while the debate continues. It is this type of connection that is added by this study.

\section{Objectives}\label{sec:objectives}

The primary purpose for the research was to eliminate excessive sycophancy through majority vote in Multi-Agent LLM Debate. The paper used a trusting mechanism instead of just using majority vote based on evidence. Therefore, the paper had to provide two initial aspects of its methodology. First, the paper needed to identify how (the) system fails. When there are multiple agents debating, a large majority (of agents) who confidently state their position as correct may cause a minority of agents who are actually correct to change their positions ~\cite{he2026minority,pitre-etal-2025-consensagent}. Second, the paper needed to develop a bounding mechanism for trust scores. Trust scores increase whenever external evidence supports a claim made by an agent and decrease whenever such evidence does not support such claims. These updates occur prior to establishing the final answer for each agent. Additionally, the paper will compare performance of the trusting mechanism to that of majority voting and other forms of debate and aggregation mechanisms ~\cite{choi2025debate,wang-etal-2024-rethinking-bounds}. Performance metrics include accuracy of answers; consensus collapse; preservation of minority; and calibration with respect to the evidence.

\section{Methodology}\label{sec:methodology}

This study employs an experimental methodology. First, we build a baseline for a debating system. Then, several diverse large language model (LLM) agents are asked the same science-based question and they revise their position based upon what they learn of each other’s reasoning through multiple rounds~\cite{du2024improving}. In addition to revising their positions based upon learning about each other’s reasoning, we have the baseline employ majority vote. This will reveal typical failures modes and allow them to be quantified~\cite{choi2025debate}. We next take the reasonings provided by the individual agents and divide those into smaller fact claims. Those fact claims are verified against scientific evidence obtained from outside data-sources, e.g., academic papers and reference databases ~\cite{karpukhin2020dense}. Any contradiction between a claim and the evidence found via retrieval would be identified and assigned a score ~\cite{kushwaha2026agentbased}. The notion that third-party tools could evaluate and/or correct model outputs provides support for the second step ~\cite{gou2024critic}. As verdicts are generated, these results will update the trust values associated with the individual agents. An approach that utilizes per-agent confidence to determine how to best direct a debate is followed here ~\cite{yoffe-etal-2025-debunc}. When a claim is supported by an agent, it increases its own trust value; when a claim is contradicted by an agent, it decreases its own trust value. Claims that cannot be definitively determined to be true or false are treated conservatively so that there is not a rapid increase/decrease in scores due to relatively weak evidence. The trust values for the individual agents remain within bounds at all times and therefore no single agent has complete control over the outcome nor does any agent ever completely lose influence. The final answer is determined via weighted aggregation of the individual agents' trust values rather than via some form of simple counting. The proposed framework was tested using scientific reasoning datasets including GPQA~\cite{rein2024gpqa} and MMLU-Pro~\cite{wang2024mmlupro}. Evaluation measures include determining if the proposed mechanism reduces consensus collapse and determining if correct minority answers are preserved. Additionally, evaluation determines if the overall accuracy of the system remains intact, again without requiring any fine tuning. Finally, Figure~\ref{fig:methodology-flow} illustrates the end-to-end process.

\begin{figure}[ht]
    \centering
    \includegraphics[width=\textwidth]{fig-pipeline.png}
    \caption{End-to-end flow of the trust-calibrated multi-agent deliberation framework: debate, claim-level evidence verification, per-round trust updates, and trust-weighted aggregation.}
    \label{fig:methodology-flow}
\end{figure}

\section{Project Outcome}\label{sec:project-outcome}

The main outcome is a reproducible framework for trust-calibrated scientific deliberation with multiple agents. Agent influence is updated from retrieved evidence while the debate runs. The final answer then rests on evidence-weighted trust, not a majority vote. The study also produces an evaluation protocol. The protocol detects sycophantic consensus~\cite{pitre-etal-2025-consensagent} and tracks how often correct minority answers survive~\cite{he2026minority}. It also measures the method against established baselines. Dynamic trust calibration should make debate more dependable on scientific questions. A wrong answer can still remain. The decision is still more transparent because the method records which agents the evidence supported and which it contradicted.

\section{Organization of the Report}\label{sec:organization-of-the-report}

The rest of the report moves from background to results. Chapter 2 covers multi-agent debate and sycophancy. It also defines the gap this study addresses. The requirements, the methodology, the system design, and the project plan sit in Chapter 3.

Chapter 4 holds the implementation setup and the testing procedure. It also presents the results with discussion. Standards, constraints, sustainability, and the complex engineering problem analysis fill Chapter 5. Chapter 6 closes with a summary and limitations. It also outlines directions for future work.
